\chapter{Implementation}
\label{ch:implementation}

%==============================================================================
\section{Traceability of Objectives to Evidence}
%==============================================================================

Software and systems requirements practice treats bidirectional traceability, from stakeholder needs through design to verification evidence, as a core quality activity \cite{ieee29148Requirements2018}.
Table~\ref{tab:obj-trace} summarises how the objectives of this thesis (hour-0 plus autoregressive ICU synthesis, mixed-type tabular handling, utility/distributional/privacy checks, reproducible logging) map to implementation anchors, evaluation quantities, and the sections below where those links are documented.

\begin{table}[htbp]
\centering
\footnotesize
\setlength{\tabcolsep}{3pt}
\begin{tabularx}{\textwidth}{@{}>{\raggedright\arraybackslash}X >{\raggedright\arraybackslash}X >{\raggedright\arraybackslash}X >{\raggedright\arraybackslash}X@{}}
\toprule
\textbf{Objective / criterion} &
\textbf{Design / implementation} &
\textbf{Metric(s)} &
\textbf{Where evidenced} \\
\midrule
Two-stage synthesis: IID hour~0 then autoregressive rows &
Numbered scripts; \code{run\_sweep} trains both stages per config &
Stage timings; end-to-end sweep rows &
section \textbf{Operational Pipeline}; section \textbf{Unified Sweep Architecture} \\
\midrule
Mixed quantitative / categorical ICU features &
\code{feature\_types} metadata; Forest-Flow vs HS3F paths; label remap in \code{sample()} &
Task inputs after \code{inverse\_transform}; per-column KS on numerics &
section \textbf{Categorical Typing}; section \textbf{Categorical Labels and Flow Integration}; section \textbf{Sampling and Reconstruction} \\
\midrule
Downstream utility of synthetic cohorts &
\code{tstr\_framework} multi-task loop &
Per-task acc., F1, ROC-AUC (macro where noted) &
section \textbf{Multi-Task TSTR}; section \textbf{Task Metrics}; Figures~\ref{fig:sweep_mort_roc}--\ref{fig:sweep_mort_f1} \\
\midrule
Distributional fidelity beyond single-task scores &
Sweep post-process: marginals, correlation gap, clinical bounds &
$\overline{\mathrm{KS}}$, $\|\Delta C\|_F$, range-violation \% &
section \textbf{Distributional Quality Metrics}; section \textbf{Empirical Outcomes from the Hyperparameter Sweep} \\
\midrule
Privacy-preservation risk checks &
\code{src/lexisflow/evaluation/privacy\_metrics.py} (DCR + membership attack) &
DCR median/match rate; baseline/overfit scores; MIA ROC-AUC &
section \textbf{Privacy Risk Metrics}; section \textbf{Key Implementation File Map} \\
\midrule
Reproducibility and auditable experiment history &
Seeds; artifact paths; canonical CSV schema; cache signatures; failure rows preserved &
Logged metrics; \code{degenerate\_flag}; resumable \code{sweep\_results.csv} &
section \textbf{Reproducibility and Reporting Conventions}; section \textbf{Unified Sweep Architecture} \\
\bottomrule
\end{tabularx}
\caption{Objective-to-evidence traceability for this implementation chapter (section titles refer to headings within this chapter).}
\label{tab:obj-trace}
\end{table}

%==============================================================================
\section{Operational Pipeline for Hyperparameter Sweep}
%==============================================================================

The current end-to-end workflow is implemented as numbered scripts with explicit artifacts:
\begin{enumerate}\tightlist
    \item \code{prepare\_hour0.py} and \code{fit\_hour0\_preprocessor.py}
          prepare and encode hour-0 training data.
    \item \code{train\_hour0.py} trains an IID generator.
    \item \code{prepare\_autoregressive.py} creates autoregressive tabular rows and carves patient-disjoint train / test / holdout splits (\code{autoregressive\_data.csv}, \code{real\_test.csv}, \code{real\_holdout.csv}).
    \item \code{fit\_autoregressive\_preprocessor.py} fits the full-data autoregressive preprocessor on the training cohort.
    \item \code{train\_autoregressive.py} trains the autoregressive generator.
    \item \code{generate.py} loads model artifacts and exports synthetic samples.
    \item \code{run\_sweep.py} performs unified hour-0 + autoregressive hyperparameter search.
\end{enumerate}

The architecture is explicitly two-stage and backbone-selectable (\code{forest-flow} or \code{hs3f}, default \code{hs3f}). Quality and TSTR evaluation are integrated directly into \code{run\_sweep.py}; ad-hoc inspection of specific generated outputs is handled by the helpers in \code{synth\_gen.evaluation}. A dedicated matched-cell driver (\code{run\_backbone\_comparison.py}) executes the same data and metric pipeline for HS3F, ForestFlow, and CTGAN at a fixed comparison cell; results are reported in the Evaluation chapter (section~\ref{sec:backbone-comparison}).
This structure was chosen so each sweep row is an auditable end-to-end experiment unit, which is why Chapter~7 can compare utility, fidelity, privacy, and runtime from one consistent log schema.

%==============================================================================
\section{Training Implementation Details}
%==============================================================================

\subsection{Hour-0 IID Training}

\code{train\_hour0.py} loads \code{models/hour0\_preprocessor.pkl}, transforms all hour-0 rows, and fits an IID generator:
\begin{equation}
\hat{p}(\mathbf{z}_0), \quad \mathbf{z}_0=[\mathbf{s},\mathbf{x}_0].
\end{equation}

The script persists model metadata including \code{model\_type}, training hyperparameters, and preprocessor linkage.
This stage exists to model initial-state realism separately from transition dynamics, which explains later findings where rollout quality and hour-0 plausibility diverge.

\subsection{Autoregressive Training}

\code{train\_autoregressive.py} loads \code{models/preprocessor\_full.pkl}, transforms selected rows of \code{autoregressive\_data.csv}, and reconstructs split matrices via persisted indices:
\begin{equation}
X_{\text{target}} = X[:, \text{target\_indices}], \quad
X_{\text{condition}} = X[:, \text{condition\_indices}].
\end{equation}

This removes the previous dependence on naive prefix slicing and is needed in the current implementation when transformed columns are type-grouped.
The design choice prevents target/condition leakage from column-order assumptions, which directly supports the validity of the reported TSTR and fidelity metrics.

\subsection{Categorical Typing Interface}

Both backbones receive \code{feature\_types} from preprocessor metadata:
\begin{itemize}\tightlist
    \item \code{'q'} quantitative,
    \item \code{'c'} categorical.
\end{itemize}

Forest-Flow internally overrides noised target dimensions to quantitative mode (flow arithmetic constraint), while HS3F keeps explicit categorical heads \cite{akazan2024hs3f}.
This was implemented to keep both backbones compatible with one data interface while preserving HS3F's mixed-type path, which is why later backbone differences are interpretable as modelling effects rather than pipeline mismatch.

\subsection{Categorical Labels and Flow Integration (Codebase)}

For HS3F categorical dimensions, training remaps observed categories to contiguous local class indices for XGBoost multiclass fitting; \code{sample()} inverts that map via stored \code{local\_to\_orig} so synthetic rows use the same encoded IDs as the preprocessor artifact.
The training remap yields contiguous class indexing under arbitrary subsampling without sparse multiclass labels at fit time.

For continuous targets under Forest-Flow, sampling follows explicit first-order integration.
Forest-Flow integrates the backward flow with explicit Euler updates (\code{X\_t += h * Y\_hat}) in \code{src/lexisflow/models/forest\_flow.py}.
HS3F continuous dimensions default to \code{solver="rk4"} in \code{src/lexisflow/models/hs3f.py}, with \code{solver="euler"} available as an alternative.

%==============================================================================
\section{Sampling and Reconstruction}
%==============================================================================

\subsection{Model-Space Sampling}

For conditional generation, models sample:
\begin{equation}
\tilde{X}_{\text{target}} \sim p_{\theta}(X_{\text{target}} \mid X_{\text{condition}}).
\end{equation}

In sweep mode, conditioning rows are sampled from the real transformed condition matrix and reused as context for synthetic targets.

\subsection{Rebuilding Full Feature Vectors}

The implementation reconstructs full transformed rows in preprocessor order:
\begin{equation}
X_{\text{full}}[:,\text{target\_indices}] \leftarrow \tilde{X}_{\text{target}}, \quad
X_{\text{full}}[:,\text{condition\_indices}] \leftarrow X_{\text{condition}}.
\end{equation}

Then \code{inverse\_transform()} restores tabular values and category labels.
This reconstruction step is required so all downstream evaluators operate in comparable tabular semantics rather than backbone-specific latent coordinates.

\subsection{Post-Decode Clinical Normalisation}

Known intervention indicators (e.g., \code{vent}, \code{vaso}) are explicitly clipped/rounded to $\{0,1\}$ before evaluation/export.
This step is intended to restore discrete $\{0,1\}$ semantics for binary treatment flags after decode.

%==============================================================================
\section{Evaluation Stack}
%==============================================================================

\subsection{Cohort-Level Data Splits}
\label{sec:cohort-splits}

All utility, fidelity, and privacy metrics are computed against real-data cohorts that are patient-disjoint from the generator's training cohort.
At preprocessing time (\code{scripts/prepare\_autoregressive.py}), the autoregressive table is partitioned at the \code{subject\_id} level into three CSVs with an 80\,/\,10\,/\,10 split drawn from a fixed RNG seed so the partition is reproducible across runs:
\begin{itemize}\tightlist
    \item \code{data/processed/autoregressive\_data.csv} -- the training cohort (80\,\% of patients), the only file the generator ever sees.
    \item \code{data/processed/real\_test.csv} -- a TSTR / fidelity reference cohort (10\,\% of patients).
    \item \code{data/processed/real\_holdout.csv} -- a privacy-evaluation holdout cohort (10\,\% of patients).
\end{itemize}

Because the split is performed on \code{subject\_id}, no patient's timesteps ever span the training cohort and an evaluation cohort.
Consequently:
\begin{itemize}\tightlist
    \item Multi-task TSTR (section \textbf{Multi-Task TSTR}) trains classifiers on synthetic data and tests them on sequences drawn from \code{real\_test.csv}; the real-trained baseline classifier is also fit and evaluated entirely within this disjoint cohort via an internal stratified split (\code{test\_size=0.3}).
    \item The distributional fidelity metrics (KS, correlation Frobenius norm, range-violation rate) are computed against a sample of \code{real\_test.csv} rather than the training matrix, so they measure distributional recovery of \emph{unseen} rows rather than reproduction of the training set.
    \item The DCR overfitting-protection and DOMIAS-style membership-inference attack use \code{real\_holdout.csv} directly as the holdout reference, giving the overfitting-vs-generalisation comparison a meaningful signal even when the generator is expressive enough to memorise.
\end{itemize}

The split is materialised once per preprocessing run and referenced explicitly by \code{AutoregressiveInputs.real\_test\_path} and \code{AutoregressiveInputs.real\_holdout\_df} in \code{src/lexisflow/sweep/data\_prep.py}, ensuring the sweep driver, the matched backbone-comparison driver, and any future evaluation scripts consume byte-identical reference cohorts.

\subsection{Multi-Task TSTR}

The active sweep uses a unified patient-level TSTR framework with two tasks:
\begin{enumerate}\tightlist
    \item Mortality classification (binary, patient-level),
    \item ICU LOS category classification (multi-class, patient-level).
\end{enumerate}

Each task uses a \code{\_PaddedSequenceClassifier}: patient trajectories are grouped by \code{subject\_id}, sorted by \code{hours\_in}, padded to a common length with a binary mask, flattened, and fed to a logistic-regression baseline. This patient-level design eliminates within-patient data leakage and captures the temporal structure the autoregressive generator is designed to preserve. The design evolution from the original row-level three-task framework and the rationale for removing vasopressor-use prediction are documented in the Evaluation chapter. The TSTR paradigm follows established synthetic health-data evaluation practice \cite{esteban2018rcgan}.

\subsection{Task Metrics}

For binary tasks, reported metrics include accuracy, F1, and ROC-AUC.
For LOS (multi-class), macro-F1 and one-vs-rest macro ROC-AUC are included.

\subsection{Distributional Quality Metrics}

In addition to task utility, sweep evaluation computes:
\begin{enumerate}\tightlist
    \item Average KS statistic across valid numeric columns,
    \item Correlation matrix Frobenius norm,
    \item Clinical range violation percentage.
\end{enumerate}

Let $\mathcal{J}$ be valid numeric columns:
\begin{equation}
\text{KS}_{\text{avg}} = \frac{1}{|\mathcal{J}|}\sum_{j\in\mathcal{J}}\sup_x |F_j^{\text{real}}(x)-F_j^{\text{synth}}(x)|.
\end{equation}
\begin{equation}
\text{CorrDiff} = \|C^{\text{real}} - C^{\text{synth}}\|_F.
\end{equation}

Columns with insufficient finite samples or zero variance are excluded before these computations to avoid undefined correlations.

\subsection{Privacy Risk Metrics}

To complement utility/fidelity metrics, the current implementation adds four privacy checks used in recent tabular-synthetic evaluation practice \cite{lautrup2024syntheval,breugel2023domias}:
\begin{enumerate}\tightlist
    \item Distance to Closest Record (DCR) summary ($\mathrm{median}$, tails, exact-match rate),
    \item DCR baseline-protection score (synthetic-vs-random closeness to real data),
    \item DCR overfitting-protection score (synthetic closeness to train vs holdout),
    \item DOMIAS-inspired membership inference attack (MIA) ROC-AUC and attacker advantage.
\end{enumerate}

For synthetic samples $\tilde{x}_i$ and real training set $\mathcal{D}_{\text{train}}$, DCR is:
\begin{equation}
\mathrm{DCR}(\tilde{x}_i,\mathcal{D}_{\text{train}})=\min_{x\in\mathcal{D}_{\text{train}}}\|\phi(\tilde{x}_i)-\phi(x)\|_2,
\end{equation}
with mixed-type embedding $\phi(\cdot)$ (min-max normalised numerics + one-hot categoricals).
Baseline protection follows the SDMetrics ratio form \cite{sdmetricsPrivacyDocs}:
\begin{equation}
\mathrm{Score}_{\text{baseline}}=\min\left(1,\frac{\mathrm{median}\,\mathrm{DCR}_{\text{synth}\rightarrow\text{train}}}{\mathrm{median}\,\mathrm{DCR}_{\text{random}\rightarrow\text{train}}}\right).
\end{equation}

Overfitting protection compares train-vs-holdout nearest-neighbour proximity for synthetic rows, where the holdout is the patient-disjoint \code{real\_holdout.csv} cohort defined in section~\ref{sec:cohort-splits}:
\begin{equation}
\mathrm{Score}_{\text{overfit}}=\min\left(1,\,2\left(1-\Pr\left[\mathrm{DCR}_{\text{synth}\rightarrow\text{train}} < \mathrm{DCR}_{\text{synth}\rightarrow\text{holdout}}\right]\right)\right).
\end{equation}

Membership risk is quantified with a DOMIAS-style local overfitting score \cite{breugel2023domias}:
\begin{equation}
s(x)=\log\frac{d_{\text{ref}}(x)+\epsilon}{d_{\text{synth}}(x)+\epsilon},
\end{equation}
where $d_{\text{synth}}$ is nearest-neighbour distance to synthetic data and $d_{\text{ref}}$ is distance to the patient-disjoint holdout cohort, so a non-zero attacker advantage can only reflect closeness to training patients rather than to the evaluation cohort itself.
Attack performance is reported via ROC-AUC (and attacker advantage) in the membership-inference paradigm \cite{shokri2017membership}.

\paragraph{Interpretation guardrail.}
DCR-style proximity metrics are useful diagnostics but should not be treated as sufficient legal/anonymity evidence on their own; recent work shows distance-only proxies can miss true membership leakage, so MIA results remain primary for privacy claims \cite{yao2025dcrdelusion}.

\subsection{Empirical Outcomes from the Hyperparameter Sweep}

The canonical log \code{results/sweep\_results.csv} stores one normalised row per completed sweep cell over the configured $(n_t, n_{\text{noise}})$ grid. The schema includes metric columns, stage timings, error fields, and optional uncertainty fields (\code{*\_std}, \code{*\_ci95}, \code{trajectory\_seed\_count}) for multi-seed runs.

This chapter documents how those values are produced and persisted. The interpretation of utility, fidelity, privacy, and efficiency outcomes is intentionally centralised in Chapter~7.

\paragraph{Trajectory-Level Coherence Metrics}
\label{sec:trajectory-coherence}

The trajectory-level metrics defined in Chapter~7 (autocorrelation distance, stay-length KS, transition smoothness ratio, temporal cross-feature drift) are implemented in \code{src/lexisflow/evaluation/trajectory\_metrics.py} and integrated into \code{run\_sweep.py}, with dedicated unit coverage in \code{src/lexisflow/evaluation/tests/test\_trajectory\_metrics.py}. Implementation correctness is established by tests that verify:
\begin{itemize}\tightlist
    \item zero values when real and synthetic trajectories are identical (sanity bound);
    \item monotonic degradation as synthetic trajectories are perturbed with progressively larger noise;
    \item stable behaviour when trajectory lengths fall below the minimum window (short trajectories are excluded rather than causing NaNs).
\end{itemize}
Numerical results and interpretation are reported in Chapter~7.

%==============================================================================
\section{Unified Sweep Architecture}
\label{sec:unified-sweep}
%==============================================================================

\subsection{Per-Configuration Lifecycle}

For each $(n_t, n_{\text{noise}})$ pair, \code{run\_sweep.py} executes:
\begin{enumerate}\tightlist
    \item Train hour-0 IID model,
    \item Train autoregressive model,
    \item Generate synthetic samples,
    \item Run multi-task TSTR + quality metrics,
    \item Append one normalised CSV row.
\end{enumerate}

Each finished configuration is normalised to the canonical header list \code{SWEEP\_RESULT\_COLUMNS} before append, so partial reruns do not corrupt \code{results/sweep\_results.csv}.
The row payload includes utility metrics, quality metrics, timing fields for both model stages, and \code{degenerate\_flag}.

\subsection{Always-On Transformed Cache}

Autoregressive transformed arrays are cached by default under \code{models/sweep\_cache/transformed\_autoregressive/} and reloaded when a signature check passes.

The signature dict in \code{metadata.pkl} is produced by \code{\_build\_cache\_signature} in \code{run\_sweep.py}: CSV path, size, and nanosecond mtime; \code{n\_cols} and \code{n\_target}; SHA-256 hashes of the ordered column-name list, \code{target\_indices}, and \code{condition\_indices}.
\code{\_load\_transformed\_cache} requires exact equality of all signature fields before reusing cached \code{.npy} arrays.

A hit skips redundant \code{transform} work; a miss rebuilds matrices and overwrites the cache, preventing reuse after the input file, schema, or index split changes.

\subsection{Resumability and Failure Logging}

Runs are keyed by \code{(nt, n\_noise)} from \code{results/sweep\_results.csv}; any existing row for a key is skipped on rerun, including rows previously written after failures.
Failures do not abort the full sweep; they are written as explicit NaN/error rows, preserving an audit trail.

\subsection{Degeneracy Guardrail}

A run is flagged as degenerate when both patient-level synthetic-side thresholds are simultaneously met:
\code{mortality\_synth\_roc\_auc <= 0.55} and \code{los\_synth\_macro\_f1 <= 0.15}, with finite-value checks before thresholding.
Vasopressor metrics are no longer part of the degeneracy check following the TSTR framework revision described in the Evaluation chapter.
Importantly, runs are \emph{not} dropped; metrics are preserved and annotated with \code{degenerate\_flag=1}.

%==============================================================================
\section{Reproducibility and Reporting Conventions}
%==============================================================================

The current report and code follow reproducibility-oriented implementation practices:
\begin{itemize}\tightlist
    \item deterministic seeds for NumPy/model constructors,
    \item explicit artifact paths and metadata-rich pickles,
    \item canonical CSV schema enforcement before append,
    \item transparent warnings instead of silent failures,
    \item progress instrumentation for long-running stages.
\end{itemize}

These choices are consistent with reproducible computational research practice and are intended to support debugging and thesis traceability.

%==============================================================================
\section{Relationship to External Methods and Projects}
%==============================================================================

Our implementation is grounded in:
\begin{itemize}\tightlist
    \item flow matching for CNF-style training \cite{lipman2023flow},
    \item tree-based flow matching for tabular generation \cite{jolicoeur2023forest},
    \item heterogeneous sequential mixed-type generation \cite{akazan2024hs3f},
    \item XGBoost as the core function approximator \cite{chen2016xgboost}.
\end{itemize}

Common synthetic tabular/time-series baselines include CTGAN \cite{xu2019ctgan}, TimeGAN \cite{yoon2019timegan}, TabDDPM \cite{kotelnikov2023tabddpm}, TabSyn \cite{zhang2024tabsyn}, and SDV ecosystem tooling \cite{sdvproject}. In this project, CTGAN is integrated through \code{src/lexisflow/models/ctgan\_adapter.py} as a pragmatic lower-bound/reference baseline inside the same data and metric pipeline. Transformer-based generators such as GReaT \cite{borisov2023great} and REaLTabFormer \cite{solatorio2023realtabformer} remain future comparison points.

%==============================================================================
\section{Key Implementation Modules}
%==============================================================================

The implementation centres on a small set of modules, each tied to a reported outcome:
\begin{itemize}\tightlist
    \item \code{src/lexisflow/data/transformers.py}: preprocessing/inverse decoding that governs whether generated rows can be evaluated in clinically meaningful units.
    \item \code{src/lexisflow/models/\{forest\_flow,hs3f\}.py}: backbone-specific training/sampling logic that drives fidelity and utility differences.
    \item \code{src/lexisflow/evaluation/\{tstr\_framework,privacy\_metrics,trajectory\_metrics\}.py}: metric implementations that produce the evidence interpreted in Chapter~7.
    \item \code{scripts/\{fit\_autoregressive\_preprocessor,train\_hour0,train\_autoregressive,run\_sweep\}.py}: orchestration layer that makes runs reproducible and sweep results comparable.
\end{itemize}
